\section*{How This Book Should Be Read}
\addcontentsline{toc}{section}{How This Book Should Be Read}

This is not a survey. A survey skims broadly over advanced topics without working through their foundations. This book focuses directly on the concrete mechanisms of proof.

We spend substantial time on the initial foundations---statements, variables, quantifiers, negations, definitions, and sets---because these ideas support everything that follows. Without a secure grasp of these tools, real analysis (the rigorous study of limits, continuity, and infinite precision) can seem obscure, abstract algebra (the study of symmetries and algebraic operations) can appear disconnected from familiar arithmetic, and topology (the study of properties preserved under continuous deformation) can seem purely abstract. 

When these foundations are clear, however, each subject develops naturally: their definitions and theorems follow directly from asking precise questions about mathematical structure.

\begin{center}
\begin{tikzpicture}[
    node distance=0.8cm,
    font=\small,
    lane/.style={draw=gray!30, fill=gray!3, rounded corners=6pt, inner sep=10pt},
    stepbox/.style={draw=#1!70!black, fill=#1!8, rounded corners=3pt, thick, align=center, inner sep=6pt, text width=2.5cm, minimum height=1.1cm},
    stepbox/.default=sourceblue,
    arrup/.style={-{Latex[length=2.5mm]}, thick, #1},
    arrup/.default=gray!60
]

% Background lanes
\node[lane, minimum width=13.5cm, minimum height=1.7cm] (lane1) at (6.0, 2.3) {};
\node[lane, minimum width=13.5cm, minimum height=1.7cm] (lane2) at (6.0, -0.4) {};

% Labels for lanes
\node[anchor=west, font=\footnotesize\scshape, text=gray!60] at (-0.5, 3.4) {Mechanical Approach: Calculation \& Pattern Guessing};
\node[anchor=west, font=\footnotesize\scshape, text=sourceblue!80!black] at (-0.5, 0.7) {Structural Approach: Pure Mathematics \& Deductive Proof};

% Top lane nodes
\node[stepbox=gray] (c1) at (1.2, 2.3) {Calculate\\Specific Cases};
\node[stepbox=gray] (c2) at (4.4, 2.3) {Observe\\Numerical Pattern};
\node[stepbox=trapred] (c3) at (7.6, 2.3) {Plausible Guess\\(Unverified)};
\node[stepbox=trapred] (c4) at (10.8, 2.3) {Vulnerable to\\Hidden Failure};

\draw[arrup=gray!60] (c1) -- (c2);
\draw[arrup=gray!60] (c2) -- (c3);
\draw[arrup=trapred!70] (c3) -- (c4);

% Transition arrows between lanes
\draw[-{Latex[length=3mm]}, very thick, quotientorange] (10.8, 1.45) -- node[right=4pt, font=\scriptsize\bfseries, text=quotientorange]{Why does it hold?} (10.8, 0.45);

% Bottom lane nodes
\node[stepbox=sourceblue] (p1) at (1.2, -0.4) {Concrete Intuition\\\& Visual Model};
\node[stepbox=structurepurple] (p2) at (4.4, -0.4) {Precise\\Definition};
\node[stepbox=structurepurple] (p3) at (7.6, -0.4) {Step-by-Step\\Deduction};
\node[stepbox=targetgreen] (p4) at (10.8, -0.4) {Verified\\General Theorem};

\draw[arrup=sourceblue!80] (p1) -- (p2);
\draw[arrup=structurepurple!80] (p2) -- (p3);
\draw[arrup=structurepurple!80] (p3) -- (p4);

\end{tikzpicture}
\end{center}

\subsection*{The anatomy of a mathematical concept}

In this text, we resist the urge to hurry into formal proofs before the picture is settled in your mind. Whenever a major concept arrives, you will encounter it through a deliberate sequence:

\begin{itemize}
\item \textbf{The geometric or tactile picture}: What does this look like if you draw it, walk through it, or treat it as a game?
\item \textbf{The minimal formal definition}: What exact requirements are packed into the vocabulary, and why is every single word necessary?
\item \textbf{Examples and deliberate non-examples}: What lives inside the concept, and what sits just outside its boundary?
\item \textbf{Notation decoded}: What is the shorthand saying in plain human terms, and what common reading traps does it set?
\item \textbf{Complete, unhurried proofs}: Every algebraic step and logical pivot explained without hand-waving or skipping the mechanics.
\item \textbf{Walking questions}: Prompts designed to be turned over in your head while away from your desk.
\end{itemize}

\begin{focusbox}
Measure your progress through this book by your ability to work with the concepts independently: after reading an explanation, close the book, take a blank sheet of paper, and verify whether you can state the definitions precisely and explain the reasoning behind the proofs.
\end{focusbox}

\subsection*{The working rule for this text}

When you encounter a new definition, test it immediately against concrete examples. Find a simple object that satisfies it, find a close non-example that fails one condition, and ask what would change if a condition were omitted.

This practice builds immediate familiarity. The theorems in advanced mathematics describe how these definitions interact, so understanding the definitions thoroughly makes following and discovering proofs much more natural.

\begin{recoverybox}
\textbf{Pause and think:} Consider the definition of an \emph{even integer}: an integer $n$ for which there exists an integer $k$ such that $n = 2k$. Test this machine on three edge cases before reading further:
\begin{itemize}
\item Is $0$ even? Check the machine: does an integer $k$ exist with $0 = 2k$? Yes, $k = 0 \in \Z$, so $0$ passes.
\item Is $-6$ even? Check the machine: $-6 = 2(-3)$, and $-3 \in \Z$, so $-6$ passes.
\item Is $2.4$ even? The definition requires $n \in \Z$; $2.4$ is not an integer, so the condition fails immediately.
\end{itemize}
Notice how a formal definition settles questions that everyday intuition might leave uncertain.
\end{recoverybox}

We can visualize a definition as a two-stage filter machine (Figure~\ref{00:fig-definition-machine}).

\begin{figure}[htbp]
\centering
\begin{tikzpicture}[
    font=\small,
    box/.style={draw=#1!70!black, fill=#1!8, rounded corners=4pt, thick, align=center, inner sep=6pt},
    box/.default=sourceblue,
    arrup/.style={-{Latex[length=2.5mm]}, thick, #1},
    arrup/.default=gray!60
]

% Input candidates
\node[box=gray, font=\footnotesize, text width=2.0cm] (in) at (0, 1.5) {\textbf{Candidates}\\$0, -6, 2.4, 7$};

% Gate 1: Domain
\node[box=sourceblue, text width=2.4cm] (g1) at (3.4, 1.5) {\textbf{Domain Gate}\\Is $n \in \Z$?};

% Rejection 1
\node[box=trapred, text width=2.3cm, font=\footnotesize] (r1) at (3.4, -0.7) {\textbf{Rejected}\\$2.4 \notin \Z$\\Not an integer};

% Gate 2: Structure
\node[box=structurepurple, text width=2.6cm] (g2) at (7.2, 1.5) {\textbf{Structure Gate}\\$\exists k \in \Z \, (n = 2k)$?};

% Rejection 2
\node[box=quotientorange, text width=2.3cm, font=\footnotesize] (r2) at (7.2, -0.7) {\textbf{Fails Condition}\\$7 = 2(3) + 1$\\Odd integer};

% Accepted
\node[box=targetgreen, text width=2.4cm, font=\footnotesize] (acc) at (10.9, 1.5) {\textbf{Accepted: Even}\\$0 = 2(0)$\\$-6 = 2(-3)$};

% Arrows
\draw[arrup=gray!70] (in) -- (g1);
\draw[arrup=sourceblue!80] (g1) -- node[above=2pt, font=\scriptsize\bfseries, text=sourceblue!80!black]{Yes} (g2);
\draw[arrup=trapred!80] (g1) -- node[right=2pt, font=\scriptsize\bfseries, text=trapred]{No} (r1);
\draw[arrup=structurepurple!80] (g2) -- node[above=2pt, font=\scriptsize\bfseries, text=structurepurple!80!black]{Yes} (acc);
\draw[arrup=quotientorange!80] (g2) -- node[right=2pt, font=\scriptsize\bfseries, text=quotientorange]{No} (r2);

\end{tikzpicture}
\caption{The definition as an inspection machine. Candidate inputs face sequential gates: first a domain check ($n \in \Z$), then the algebraic condition ($\exists k \in \Z$ with $n = 2k$). Candidates that satisfy both are accepted as even; failing either halts the process.}
\label{00:fig-definition-machine}
\end{figure}

\subsection*{From measurement to deductive proof}

Using empirical rules and arithmetic calculations, Egyptian and Babylonian scribes computed canal volumes, granary capacities, and calendar cycles with notable practical accuracy.

Beginning around the sixth century BCE, Greek mathematicians pursued a different goal: explaining why geometric relationships must hold in general. Consider an isosceles triangle (a triangle with two sides of equal length). A surveyor can measure its base angles with a protractor, verify their equality across many drawings, and observe that the base angles match.

Why, then, did Euclid construct a formal \emph{proof}?

Because a protractor inspects a single physical drawing. It is subject to eyesight and the limits of the instrument. It can confirm that the base angles match within a fraction of a degree, but it cannot establish that the difference is strictly zero. More fundamentally, measuring a thousand triangles cannot guarantee that every possible isosceles triangle shares the same property.

Deductive proof shifts the focus from measurement to structure (Figure~\ref{00:fig-measurement-vs-proof}). Instead of measuring the angles, we draw the vertical axis of symmetry $AD$. By Side-Angle-Side congruence (triangles sharing two sides and the angle between them are identical), the left half $\triangle ABD$ and the right half $\triangle ACD$ are mirror reflections. The equality $\angle B = \angle C$ does not depend on measurement precision; it follows deductively from the symmetry of the figure.

\begin{figure}[htbp]
\centering
\begin{tikzpicture}[
    font=\small,
    panel/.style={draw=gray!30, fill=gray!2, rounded corners=6pt, inner sep=8pt}
]

% Panel 1: Empirical Measurement
\node[panel, minimum width=6.3cm, minimum height=5.6cm] (p1) at (3.15, 2.0) {};
\node[font=\footnotesize\scshape, text=quotientorange!90!black] at (3.15, 4.45) {Empirical Measurement (Protractor)};

% Triangle in Panel 1
\coordinate (A1) at (3.15, 3.8);
\coordinate (B1) at (1.2, 2.1);
\coordinate (C1) at (5.1, 2.1);

\draw[thick, draw=gray!70, fill=gray!8] (A1) -- (B1) -- (C1) -- cycle;

% Dimension labels
\node[left=2pt, font=\scriptsize, text=gray!80!black] at (2.0, 3.0) {side $a$};
\node[right=2pt, font=\scriptsize, text=gray!80!black] at (4.3, 3.0) {side $a$};
\node[below=2pt, font=\scriptsize, text=gray!80!black] at (3.15, 2.1) {base};

% Angle arcs and callouts
\draw[thick, quotientorange] (1.7, 2.1) arc (0:47:0.5);
\draw[thick, quotientorange] (4.6, 2.1) arc (180:133:0.5);

\node[font=\tiny\bfseries, text=quotientorange!95!black, align=center] at (1.9, 1.45) {$\angle B \approx 47^\circ \pm 0.5^\circ$};
\node[font=\tiny\bfseries, text=quotientorange!95!black, align=center] at (4.4, 1.45) {$\angle C \approx 47^\circ \pm 0.5^\circ$};

\node[draw=trapred!60, fill=trapred!8, rounded corners=3pt, text width=5.6cm, align=center, inner sep=4pt] at (3.15, 0.2) {
    \scriptsize\textbf{Empirical Limit:}\\
    Instrument precision is finite ($\pm \epsilon$). Checking 1,000 stone monuments cannot rule out a discrepancy of $0.0001^\circ$ or verify triangles yet unbuilt.
};

% Panel 2: Deductive Symmetry
\node[panel, minimum width=6.3cm, minimum height=5.6cm] (p2) at (10.05, 2.0) {};
\node[font=\footnotesize\scshape, text=sourceblue!90!black] at (10.05, 4.45) {Deductive Proof (Structural Symmetry)};

% Triangle in Panel 2
\coordinate (A2) at (10.05, 3.8);
\coordinate (B2) at (8.1, 2.1);
\coordinate (C2) at (12.0, 2.1);
\coordinate (D2) at (10.05, 2.1);

% Shaded congruent halves
\draw[fill=sourceblue!15, draw=sourceblue!80, thick] (A2) -- (B2) -- (D2) -- cycle;
\draw[fill=targetgreen!15, draw=targetgreen!80, thick] (A2) -- (D2) -- (C2) -- cycle;

% Symmetry spine
\draw[very thick, dashed, structurepurple!80] (A2) -- (D2) node[midway, right=1pt, font=\tiny\bfseries, text=structurepurple!90!black] {Axis};

% Labels
\node[above=2pt, font=\scriptsize] at (A2) {$A$};
\node[left=2pt, font=\scriptsize] at (B2) {$B$};
\node[right=2pt, font=\scriptsize] at (C2) {$C$};
\node[below=2pt, font=\scriptsize] at (D2) {$D$};

% Angle marks
\draw[thick, sourceblue!90!black] (8.6, 2.1) arc (0:47:0.5);
\node[font=\tiny\bfseries, text=sourceblue!90!black] at (8.9, 2.35) {$\angle B$};

\draw[thick, targetgreen!90!black] (11.5, 2.1) arc (180:133:0.5);
\node[font=\tiny\bfseries, text=targetgreen!90!black] at (11.2, 2.35) {$\angle C$};

% Congruence note
\node[draw=targetgreen!70, fill=targetgreen!8, rounded corners=3pt, text width=5.6cm, align=center, inner sep=4pt] at (10.05, 0.2) {
    \scriptsize\textbf{Structural Guarantee:}\\
    Angle bisector $AD$ creates $\triangle ABD \cong \triangle ACD$. Deductive congruence ensures $\angle B = \angle C$ exactly across all isosceles triangles.
};

\end{tikzpicture}
\caption{Measurement versus proof. \textbf{Left}: A physical protractor inspects a concrete triangle. It confirms equality up to instrument tolerance ($\pm \epsilon$), but cannot guarantee zero discrepancy or speak about triangles not yet measured. \textbf{Right}: A deductive argument constructs the axis of symmetry $AD$. By Side-Angle-Side congruence, the left half is an exact mirror reflection of the right half, establishing that the base angles are equal for every isosceles triangle.}
\label{00:fig-measurement-vs-proof}
\end{figure}

\begin{historybox}
The limits of physical measurement became sharp with the Pythagorean discovery of incommensurability (quantities that share no common measuring unit). The early Pythagoreans believed every geometric relationship was a ratio of whole numbers. When measuring the diagonal of a unit square against its side, they expected a fine enough unit to divide evenly into both. Instead, deductive proof revealed that no such unit exists: the ratio is $\sqrt{2}$, an irrational number (it cannot be written as a fraction of integers $a/b$). Legend says Hippasus of Metapontum was drowned at sea for revealing this result. Physical measurement produces approximations like $1.414$ or $1.41421356$, but only deductive proof could show that the ratio is not a fraction.

Euclid's breakthrough in the \emph{Elements} (c.~300 BCE) was to replace physical measurement with an axiomatic framework. Starting from minimal definitions and five geometric postulates (such as ``a straight line can be drawn between any two points''), each proposition was deduced step by step.

Even so, Euclid's system left gaps. His synthetic geometry (geometry built directly from axioms without coordinates) relied implicitly on diagrams---for instance, assuming without proof that two circles must intersect if their centers are close enough. In 1899, David Hilbert supplied a complete set of axioms. In the nineteenth century, mathematicians also discovered that Euclid's Fifth Postulate (the parallel postulate) can be altered, yielding consistent non-Euclidean geometries. In 1931, Kurt G\"odel showed that no single axiomatic system can be both consistent and complete enough to prove all truths of arithmetic. The search for secure foundations remains an active part of mathematical thinking.
\end{historybox}

The proof that the diagonal of a unit square cannot be expressed as a ratio of whole numbers illustrates the difference between physical approximation and deduction (Figure~\ref{00:fig-incommensurability-diagonal}).

\begin{figure}[htbp]
\centering
\begin{tikzpicture}[
    font=\small,
    box/.style={draw=#1!70!black, fill=#1!8, rounded corners=4pt, thick, align=left, inner sep=6pt},
    box/.default=sourceblue
]

% Left: The unit square
\draw[thick, fill=gray!5] (0, 0) rectangle (3.2, 3.2);
\draw[very thick, sourceblue] (0, 0) -- (3.2, 3.2);

% Right angle mark at (3.2, 0)
\draw[gray!60] (2.9, 0) -- (2.9, 0.3) -- (3.2, 0.3);

% Labels on the square
\node[below=2pt, font=\small] at (1.6, 0) {side $1$};
\node[right=2pt, font=\small] at (3.2, 1.6) {side $1$};
\node[above left=-1pt, font=\small\bfseries, text=sourceblue!90!black] at (1.4, 1.8) {$d = \sqrt{2}$};

% Caliper approximation callout below
\node[font=\footnotesize, text=quotientorange!90!black, align=center] at (1.6, -0.75) {Calipers: $d \approx 1.4142\dots$};

% Right: Contradiction steps
\node[box=sourceblue, text width=7.6cm] (step1) at (9.0, 2.5) {
    \textbf{1. Assume rational common unit:}\\
    $\sqrt{2} = \frac{p}{q}$ with integers $p, q \in \Z$ and $\gcd(p,q) = 1$.
};

\node[box=structurepurple, text width=7.6cm] (step2) at (9.0, 1.1) {
    \textbf{2. Square both sides:}\\
    $2 = \frac{p^2}{q^2} \implies p^2 = 2q^2 \implies p$ is even ($p = 2m$).
};

\node[box=trapred, text width=7.6cm] (step3) at (9.0, -0.45) {
    \textbf{3. Deductive contradiction:}\\
    $2q^2 = (2m)^2 = 4m^2 \implies q^2 = 2m^2 \implies q$ is even.\\
    Both $p$ and $q$ are even, contradicting $\gcd(p,q) = 1$!
};

\draw[-{Latex}, thick, gray!60] (step1) -- (step2);
\draw[-{Latex}, thick, gray!60] (step2) -- (step3);

\end{tikzpicture}
\caption{Incommensurability of the diagonal. Physical calipers can only approximate $d \approx 1.4142$, never ruling out a rational fraction with a large denominator. Deductive proof establishes that assuming $\sqrt{2} = p/q$ forces both $p$ and $q$ to be even, an impossible contradiction that proves no such fraction exists.}
\label{00:fig-incommensurability-diagonal}
\end{figure}

\subsection*{What beginner-friendly means here}

Being accessible does not mean avoiding technical precision; it means ensuring every logical step is clearly explained. We introduce concepts with concrete motivation and visual intuition, but every explanation leads directly to precise definitions and verifiable deductions. The goal is to make rigorous mathematical reasoning understandable from the ground up.

In Chapter~1 (\emph{The Change in Game: From Getting Answers to Giving Reasons}), we make this transition concrete. We will examine a formula that holds for billions of consecutive numbers before abruptly failing, and see why numerical calculation cannot substitute for deductive proof.

\subsection*{Exercises}

\begin{enumerate}
\item \textbf{Laser measurement versus deductive proof.} Suppose you build a physical triangle out of machined steel rods whose lengths agree to within one micrometer ($10^{-6}\text{ m}$). A laser sensor measures the three angles and reports $59.999^\circ$, $60.001^\circ$, and $60.000^\circ$. Have you proven that an equilateral triangle has equal angles? What does a deductive proof provide that the physical sensor can never deliver?

\item \textbf{The machine test: why $1$ is not prime.} A student argues: ``Since $1$ is only divisible by $1$ and itself, $1$ must be a prime number.'' Mathematicians define a prime number as an integer $p > 1$ whose only positive divisors are $1$ and $p$. Why is the requirement $p > 1$ deliberately included in the definition machine? (Hint: Consider what happens to the Fundamental Theorem of Arithmetic---which states that every integer greater than $1$ factors into prime numbers in exactly one way---if $1$ were counted as a prime.)

\item \textbf{Boundary testing on a definition machine.} Consider the definition: ``An integer $n$ is \emph{composite} if there exist integers $a > 1$ and $b > 1$ such that $n = ab$.'' Test each candidate number against this machine and state whether it is accepted, rejected, or outside the domain of definition:
\begin{enumerate}[label=(\alph*)]
\item $n = 12$
\item $n = 7$
\item $n = 1$
\item $n = 0$
\item $n = -6$
\end{enumerate}

\item \textbf{Testing claims on numbers.} Consider the claim: ``For every real number $x$, if $x > 0$, then $x^2 > x$.''
\begin{enumerate}[label=(\alph*)]
\item Find two numbers that easily satisfy this claim.
\item Test the boundary value $x = 1$. Does $1^2 > 1$?
\item Test the value $x = 1/2$. Does $(1/2)^2 > 1/2$?
\item Describe the exact collection of positive real numbers for which the claim $x^2 > x$ is actually true.
\end{enumerate}

\item \textbf{Incommensurability and proof.} In your own words, explain why the discovery that the diagonal of a unit square cannot be written as a fraction of integers ($a/b$) could never have been made using bronze calipers or numerical tables alone. What role did deductive contradiction play in uncovering this fact?
\end{enumerate}

\begin{walkbox}
Consider a mathematical formula you frequently use, such as the quadratic formula or the Pythagorean theorem. Can you derive it from first principles, or do you rely on it as a memorized procedure? What definitions and intermediate steps are required to turn that formula into a complete deductive proof?
\end{walkbox}

\newpage
